\section{Transferring the Gaffney domain}\label{sec:comparison}

The domain of $d$ depends only on the quasi-isometry class of a
metric; that of $\delta$ does not. We now give the additional
estimates which permit the cone theorem to be transferred to the
moduli space. All statements about domains in an open set concern
forms compactly supported in that set.

The distinction can be seen from the formulas for the operators.
The weak exterior derivative is independent of the metric, and
uniformly equivalent metrics give equivalent $L^2$ norms on forms
and their exterior derivatives. The codifferential involves the
Hodge star and therefore derivatives of the metric. Our comparisons
will control those derivatives and require the relevant metric error
to vanish in the directions where derivatives of a form may grow.
The weighted estimate below specifies the amount of vanishing that
is needed.

\subsection{Weighted interior regularity}
Put $\rho(x)=\operatorname{dist}(x,X_{\sing})$. We require the
following elementary form of bounded geometry: there are charts on
balls of radius comparable to $\rho(x)$ in which the rescaled
metric $\rho(x)^{-2}g$ is uniformly elliptic with uniformly bounded
first derivatives, and the measure is locally doubling, with uniform
constants. More precisely, charts from the Euclidean unit ball
contain $\mathbb B(x,\kappa\rho(x))$ in their half-ball image and
have image in $\mathbb B(x,\rho(x)/2)$, for fixed $\kappa>0$.
The bounds are understood on each region under discussion.
Products of a compact smooth manifold with a smooth-link cone have
this property.

\begin{lemma}\label{lem:weighted}
Under these bounds, forms in $V$ with support in a fixed compact
set satisfy
\begin{equation}\label{eq:weighted}
 \|\rho\nabla u\|\le C(\|\rho Du\|+\|u\|).
\end{equation}
\end{lemma}
\begin{proof}
Interior elliptic regularity first gives $u\in W^{1,2}_{\rm loc}$
on the regular locus. In a rescaled chart the principal coefficients
of $d+\delta$ are uniformly elliptic and Lipschitz, and its
zeroth-order coefficients are bounded. Freezing the principal
coefficients on sufficiently small balls and using the Fourier
estimate gives a uniform interior G\aa rding inequality. Rescaling
back gives
\[
 \rho(x)^2\int_{\mathbb B(x,\kappa\rho(x))}|\nabla u|^2
 \le C\int_{\mathrm{chart}}(\rho(x)^2|Du|^2+|u|^2).
\]
Since $\rho$ is Lipschitz, its values are comparable on such a
chart. A maximal disjoint family of smaller balls has enlarged
balls of bounded overlap by doubling. Summation proves
\eqref{eq:weighted}.
\end{proof}

\begin{definition}\label{def:adapted}
A chart from a model $(X_0,g_0,I_0)$ to $(X,g,I)$ is
\emph{adapted} if it is a homeomorphism, a diffeomorphism on regular
loci, and its pullback $(g',I')$ satisfies
\begin{equation}\label{eq:adapted}
 cg_0\le g'\le c^{-1}g_0,\quad
 |g'-g_0|+|I'-I_0|\le C\rho,\quad
 |\nabla^0g'|+|\nabla^0I'|\le C.
\end{equation}
The model is K\"ahler on its regular part; all norms use $g_0$.
An adapted metric chart requires only the metric bounds.
\end{definition}

\begin{theorem}\label{thm:comparison}
Suppose the model has the bounded geometry above. An adapted
metric chart identifies its Gaffney domain with that of $g'$, with
equivalent graph norms on a fixed compact support. If the chart is
adapted also for the complex structures, Weil invariance transfers
from the model to $g'$.
\end{theorem}
\begin{proof}
The same rescaled charts give bounded geometry for $g'$, and the
two connections differ by a bounded tensor. Hence
\eqref{eq:weighted} holds for both metrics with $\nabla^0$.
Let $S=*_{0}^{-1}*'$ on the graded exterior bundle. Then
\[
 \delta'=S^{-1}\delta_0S,\qquad |S-1|\le C\rho,
 \qquad |\nabla^0S|\le C.
\]
For a local orthonormal frame $e_i$, differentiation gives
\begin{equation}\label{eq:star-commutator}
 \delta_0(Su)-S\delta_0u
 =-\sum_i\iota_{e_i}(\nabla^0_{e_i}S)u
   -\sum_i[\iota_{e_i},S]\nabla^0_{e_i}u.
\end{equation}
The right side is bounded by $C(|u|+\rho|\nabla^0u|)$.
The last sum is a first-order term and is essential. Estimate
\eqref{eq:weighted} controls it from either graph norm, proving
both domain inclusions and their norm estimates.

For the complex structures put $Q=W_0^{-1}W'$. The Weil operator
on degree $j$ is the $j$th exterior power of the complex structure
on covectors. Thus $|Q-1|\le C\rho$ and $|\nabla^0Q|\le C$.
The Leibniz formulas for $d(Qu)$ and $\delta_0(Qu)$ contain the
same bounded terms and commutators of size $O(\rho)$ multiplying
$\nabla^0u$. They show $Q V(g_0)\subset V(g_0)$. Weil invariance
of the model then gives $W'u=W_0Qu\in V(g_0)=V(g')$.
\end{proof}

The vanishing in \eqref{eq:adapted} cannot be omitted. On the flat
cone of angle $4\pi$, use $z=w^2$ and $g_0=|dz|^2$.
A cutoff of $dw$ is closed and coclosed near zero and belongs to
$V(g_0)$. For $g_b=|dz|^2+\operatorname{Re}(b\,dz^2)$,
$0<|b|<1$, the metric difference is parallel, but
$\delta_{g_b}dw$ has homogeneity $-3/2$ in $z$ and is not $L^2$.
Thus even a parallel metric perturbation can change the domain.

\subsection{Product estimates}
\begin{lemma}\label{thm:product}
Let $B$ be a closed Riemannian manifold and let $Y$ satisfy the
Stokes property. Then $B\times Y$ satisfies that property and
\begin{equation}\label{eq:tangent}
 \|\nabla_Bu\|\le C(\|Du\|+\|u\|).
\end{equation}
If both factors are K\"ahler on their regular loci and $Y$ has
Weil invariance locally, so does the product. If $Y=C(Z)$ with
admissible link, the graph-norm unit ball on a fixed compact
support is precompact in $L^2(B\times C(Z))$.
\end{lemma}
\begin{proof}
Let $\mathcal F_\lambda$ be the finite-dimensional eigenspaces of
the Hodge Laplacian on $B$. They are preserved by $D_B$ and the
degree sign $s_B$. On the product,
$D=D_B\otimes1+s_B\otimes D_Y$. Testing weak derivatives against
an eigenform on $B$ shows that each projection $u_\lambda$ lies
in $\mathcal F_\lambda\otimes V(Y)$, and
\begin{equation}\label{eq:product-energy}
 \|Du_\lambda\|^2=\lambda\|u_\lambda\|^2
       +\|(1\otimes D_Y)u_\lambda\|^2.
\end{equation}
The cross terms cancel because $D_Bs_B=-s_BD_B$ and $D_Y$ is
self-adjoint. For Stokes, take arbitrary $v\in\Dom d_{\max}$
and $w\in\Dom\delta_{\max}$, rather than assuming either lies
in $V$. Project both onto a finite sum of the spaces
$\mathcal F_\lambda$. The projections commute with $d_B$
and $\delta_B$. Their coefficient equations place the
coefficients of $v$ in $\Dom d_{Y,\max}$ and those of $w$
in $\Dom\delta_{Y,\max}$. Stokes on $Y$ and integration by
parts on $B$ give the adjoint identity for the projections.
The projections of $v,dv,w,\delta w$ converge in $L^2$,
so the identity holds on the full maximal domains.
Summing \eqref{eq:product-energy} and applying the elliptic
estimate on $B$ proves \eqref{eq:tangent}.

If $B$ is K\"ahler, $W_B$ preserves each $\mathcal F_\lambda$.
Use \eqref{eq:WI-estimate} for $Y$ on a fixed compact support in
\eqref{eq:product-energy}. The resulting estimates for
$(W_B\otimes W_Y)u_\lambda$ sum, proving Weil invariance.
For compactness, the terms of large $\lambda$ have uniformly
small norm by \eqref{eq:product-energy}; finitely many remaining
terms are compact by Lemma \ref{lem:compact-cone}.
\end{proof}

These statements apply locally when the smooth factor is not
compact. For the K\"ahler assertion one can extend a small
coordinate ball to a compact K\"ahler manifold: write its
potential as $|z|^2+O(|z|^3)$, cut off the cubic remainder at
radius $\epsilon$, and place it in a flat torus. The change in
the Levi form is $O(\epsilon)$ and preserves positivity. For
metric estimates alone a cutoff extension of the metric suffices.

\subsection{Comparison at two strata}
At the deepest stratum there are two geometrically different small
quantities. The cone radius $r$ measures distance to the vertex
stratum $B\times\{0\}$. The distance $\rho$ measures proximity
to the full singular set, including the edge of the cone. Away
from that edge the two are comparable. Near an edge point at radius
$r_y>0$, however, $\rho$ tends to zero while $r$ remains comparable
to $r_y$. An error bounded by $Cr$ need not be bounded by
$C\rho$ there.

This is precisely where the product estimate is used. The full
gradient is controlled only after multiplication by $\rho$;
the derivative tangent to the smooth edge factor is controlled
without that vanishing weight. A perturbation pulled back from
the smooth factor may consequently have size $O(r_y)$, provided
the remainder is $O(\rho)$. In the Hodge-star commutator, the
smooth-factor perturbation commutes with contraction in normal
directions and therefore multiplies only tangential derivatives.
The two estimates are arranged as follows:
\begin{center}
\begin{tabular}{lll}
\toprule
Region or direction & Permitted coefficient & Controlled derivative\\
\midrule
Away from the edge, $\rho\asymp r$ & $O(r)=O(\rho)$ & $\rho\nabla u$\\
Along the smooth edge factor & $O(r_y)$ & $r_y\nabla_{B_y}u$\\
Remainder near the edge & $O(\rho)$ & $\rho\nabla u$\\
\bottomrule
\end{tabular}
\end{center}
The following definition expresses this directional distinction
in coordinates that remain uniform as the edge approaches the
vertex stratum.

Let $X_0=B\times C(Z)$ have an edge away from its vertex stratum.
Assume bounded geometry at scale $\rho$ and Weil invariance on
$X_0$. At each edge point $y$, put $r_y=r(y)$. We require
\emph{uniform edge charts} with product model
$P_y=B_y\times C(Z')$, where $Z'$ is a fixed smooth closed
odd-dimensional link. Their images contain
$\mathbb B(y,\kappa r_y)$ and lie in
$\mathbb B(y,r_y/2)$. Their domains are closed under radial
contraction of the cone factor. After rescaling by $r_y^{-1}$,
the smooth factors have uniform first-order geometry and the
charts satisfy \eqref{eq:adapted} with uniform constants. For
the complex structures the smooth factors need only carry
uniformly controlled almost complex structures. In these charts
$\rho$ is comparable to the cone radius near $y$.

\begin{definition}\label{def:two-level}
A chart at the vertex stratum is \emph{adapted to both strata}
if it is a homeomorphism and regular-locus diffeomorphism, and
its pullback tensors satisfy
\begin{equation}\label{eq:two-level}
 cg_0\le g'\le c^{-1}g_0,\quad
 |g'-g_0|+|I'-I_0|\le Cr,\quad
 |\nabla^0g'|+|\nabla^0I'|\le C,
\end{equation}
and, in every specified edge chart lying in its domain,
\begin{equation}\label{eq:edge-split}
 g'-g_0=t_y+e_y,\qquad I'-I_0=t'_y+e'_y.
\end{equation}
Here $t_y,t'_y$ are tensors pulled back from $B_y$,
$|t_y|+|t'_y|\le Cr_y$, $|\nabla t_y|+|\nabla t'_y|\le C$,
$|t_y|<1/2$, and $|e_y|+|e'_y|\le C\rho$.
\end{definition}

\begin{theorem}\label{thm:two-level}
Under these assumptions the chart identifies the Gaffney domains
and transfers Weil invariance on a smaller neighbourhood of the
vertex stratum.
\end{theorem}
\begin{proof}
Shrink to the points $x$ with $r(x)<r_*/2$ and
$\mathbb B(x,2r(x))$ contained in the chart domain, where $r_*$
is the edge-chart radius bound. This still contains a neighbourhood
of each point of the vertex stratum. Bounded geometry and
\eqref{eq:weighted} hold for $g_0$ and $g'$ by
\eqref{eq:two-level}.

For a form $w$ supported in a ball of radius $c r_y$ in an edge
chart, Theorem \ref{thm:comparison} and
\eqref{eq:tangent}, rescaled, give
\begin{equation}\label{eq:edge-tangent}
 r_y\|\nabla_{B_y}w\|\le
 C(r_y\|D_0w\|+\|w\|).
\end{equation}
The same bound holds with $D'$: compare $g'$ with the product
metric $(g_{B_y}+t_y)\oplus g_{C(Z')}$.
Their difference is $O(\rho/r_y)$ with derivative $O(1/r_y)$.
This follows from \eqref{eq:two-level}, the derivative bounds
on $t_y$, and the bounds in the model edge chart; it does not
require an additional derivative assumption on $e_y$.

Again let $S=*_{0}^{-1}*'$. Its derivative is bounded and
$|S-1|\le Cr$. Where $\rho\ge c r$, the first-order sum in
\eqref{eq:star-commutator} is controlled by
$C\rho|\nabla^0u|$. In an edge chart, smooth dependence of the
Hodge star on its metric and \eqref{eq:edge-split} give
\begin{equation}\label{eq:base-star}
 S=S_{B_y}\otimes1+O(\rho),\qquad |S_{B_y}-1|\le Cr_y.
\end{equation}
To check the remainder, compare the pair of metrics with
$(g_{P_y},g_{P_y}+t_y)$. The error is bounded by
$|e_y|+|g_0-g_{P_y}||t_y|\le C\rho$; the second product is
$O((\rho/r_y)r_y)$. Contraction in a normal direction commutes
with $S_{B_y}\otimes1$. The first-order sum is consequently
bounded pointwise by
\begin{equation}\label{eq:edge-error}
 C\big(r_y|\nabla_{B_y}u|+\rho|\nabla^0u|+|u|\big).
\end{equation}
Changing connections costs $O(1/r_y)$, multiplied by $O(r_y)$,
which explains the last term.

Cover $\{\rho<cr\}$ by edge balls of radii comparable to $r_y$
with bounded overlap, using doubling. Their cutoffs have derivative
$O(1/r_y)$. Apply \eqref{eq:edge-tangent} to the cutoffs of $u$
and sum their squared estimates; cutoff errors are bounded by
$C\|u\|^2$. Equation \eqref{eq:weighted} controls the normal
term in \eqref{eq:edge-error}. Thus the right side of
\eqref{eq:star-commutator} is $L^2$, bounded by the graph norm
of either $D_0$ or $D'$. This proves both domain inclusions.
The choice of smaller neighbourhood ensures that the charts used
in this cover lie inside the original chart domain: a nearest
singular point to $x$ with $\rho(x)<cr(x)$ is an edge point $y$
with $r_y\asymp r(x)$.

Finally $Q=W_0^{-1}W'$ has the same properties as $S$ and the
same decomposition \eqref{eq:base-star}. In a product chart the
Weil operator is a tensor product, and its dependence on $I$ is
polynomial. Applying the preceding estimates to its commutators
with wedge product and contraction gives $QV(g_0)\subset V(g_0)$.
The last step of Theorem \ref{thm:comparison} transfers Weil
invariance.
\end{proof}
